\chapter{Анализ предметной области и существующих решений}
\section{Основные определения и постановка задачи}
Пусть доступны исходные домены $\mathcal{E}_s$ с распределениями $P_e(X,Y)$. Обучение и выбор гиперпараметров используют только эти домены. Для нового домена $t\notin\mathcal{E}_s$ исследуется риск
\begin{equation}
R_t(f)=\mathbb{E}_{(X,Y)\sim P_t}\ell(f(X),Y).
\end{equation}
Наблюдения целевого домена не участвуют в выборе модели. Качество source-валидации и качество переноса рассматриваются как различные характеристики процедуры обучения.

\section{Типы распределительного сдвига}
В эксперименте различаются изменения распределения признаков, условного механизма ответа, долей классов, шума, пропусков и ложных корреляций. Сценарий с пропусками характеризует изменение измерений; он не объявляется чистым covariate shift. При изменении условного механизма существование общей инвариантной функции нельзя предполагать без дополнительных условий.

\section{Рассматриваемые методы}
IRF изменяет критерий построения деревьев с учётом неоднородности доменов~\cite{irf}. CRF использует ограничения, связывающие качество классификации и доменную информацию~\cite{crf}. Meta-forests объединяет леса, обученные по исходным доменам~\cite{meta}.

MaxRM рассматривает минимизацию наибольшего эмпирического риска~\cite{maxrm}. Era splitting меняет критерий разбиения с использованием групп наблюдений~\cite{era}. Anchor boosting использует анкерную постановку для переноса модели~\cite{anchor}.

\section{Реализации и критерии сравнения}
IRF, CRF и MaxRM в стенде являются независимыми реализациями; Meta-forests --- явно обозначенной интерпретацией. Они не называются точными репликациями авторских исследований. Отдельно учитываются capacity, допустимые разбиения, bootstrap, политика признаков и source-only настройка. В качестве сильных baseline рассматриваются Random Forest, ExtraTrees и HistGBDT.

\section{Вывод по аналитической части}
Сравнение требует раздельно оценивать абсолютное качество модели, преимущество над собственным контролем и воспроизводимость. Совпадение только числа деревьев недостаточно для изоляции механизма.
